\chapter{Discussion}


\section{Validation of Result : Cluster Cepheid}

In the dataset, 17 Cepheids are identified as members of clusters with statistical parallaxes measured by \textcite{cruz2023}. A comparison between the calibrated distances of these Cepheids shows a good agreement in reddening value with $\sigma = 0.049$. 

For distance modulus, nine of 17 cluster Cepheid deviated by 0.1 mag only, but five Cepheids (EV Sct, QZ Nor, SV Vul, Y Sct and X Lac) have deviation larger than 0.4 mag.



\begin{figure}[!ht]
    \centering
    \begin{subfigure}[b]{0.49\textwidth}
        \centering
        \includegraphics[width=\linewidth]{./plots/9_compare/jesper_cruz_mu_g.pdf}
        \caption{Distance correction comparison (Gaia)}
    \end{subfigure}
    \hfill
    \begin{subfigure}[b]{0.49\textwidth}
        \centering
        \includegraphics[width=\linewidth]{./plots/9_compare/jesper_cruz_mu.pdf}
        \caption{Distance modulus comparison (IRSB)}
    \end{subfigure}
        \caption{\small Deviation in corrected distance with 17 Cluster Cepheids (Cruz, 2023)}
    \label{dis_correction}
\end{figure}




metallicity is wrong then color is too

\section{Reddening Comparison}
\cite{2009ApJ...696.1498M}


\section{Distance Comparison}

 
\section{Luminosity Comparison}

\begin{figure}[h]
	\centering
	\includegraphics[width=\linewidth]{M_V}
\end{figure}

\subsubsection*{Effect of Metallicity}
While metallicity can affect the zero-point and slope of the period--luminosity relation, it is not explicitly considered in this thesis. However, the framework developed here is readily extendable to incorporate metallicity corrections in future studies, enabling a more comprehensive calibration of the Leavitt Law across different galactic environments.


\subsubsection*{Limitation of the method}

The method primarily does not calibrate the Leavitt Law itself, rather correct distance and reddening of individual Cepheids according to the line of linear regression for given set of Cepheid. Therefore, for different dataset, correction in distance and reddening could be slightly different due to the selection bais. To minimize such effect, I had ommitted such Cepheids from the analysis which cause significant bais to regression line of residual correlation.